\section{Jelenlegi állapot és trendek}
A mai webes technológiai stackek gyakran mikro­szolgáltatásokat (API-kon keresztül kommunikáló, elkülönült szolgáltatásokat) használnak, amelyeket sokszor szerver nélküli környezetben (serverless) futtatnak, például az AWS Lambda szolgáltatásával, amelyet 2014-ben vezettek be \cite{aws2014lambda}.

A biztonság és az adatvédelem kiemelten fontossá vált: a HTTPS mára szinte mindenhol elterjedt \cite{httpArchive2025security}. A Let's Encrypt automatizálta a tanúsítványok kiadását, és a böngészők egyre inkább mindenhol titkosítást követelnek meg \cite{letsencrypt2015launch,googleChrome2018notSecure}. A tartalomszolgáltató hálózatok (CDN-ek) és az edge computing javítják a teljesítményt \cite{cloudflareCdnEdgeServer}.

A fejlesztési folyamatok iteratívvá váltak: az Agile/DevOps csapatok gyakori frissítéseket telepítenek, és automatizált tesztelési eszközöket valamint monitorozási megoldásokat használnak \cite{agileManifesto2001,dora2018accelerate}. A webes szabványok folyamatosan fejlődnek, miközben a WHATWG és a W3C a HTML, az Ecma International pedig a JavaScript alapját adó ECMAScript továbbfejlesztésén dolgozik \cite{whatwgW3c2019Mou,ecma2024}. A böngészők új technológiákkal (például WebAssembly) bővítik a platformot \cite{w3c2019wasm}.

A webes és a natív alkalmazások közötti határ egyre inkább elmosódik, ahogy a PWA-k egyre népszerűbbé válnak \cite{mdnWhatIsPwa}. A hibrid keretrendszerek és platformok, mint a React Native vagy a Flutter web, szintén ezt a közeledést erősítik \cite{reactNativeDocs,flutterWebsite}.
